\section{Life, Mind, Consciousness, and Freedom}

The tower of selves does not stop at bodies. Read at higher rungs the same construction passes through life, mind, intelligence, and a graded sense of consciousness, and it ends at the old question of freedom. None of this adds to the base. Each is the eight elements seen from further out, and the honest status of each, from solid to frankly speculative, is kept in view.

Life is what the lean looks like at the scale of a self that must hold its own order. Charge is conserved, so order drains, and a pattern that is to persist against the drain must actively pump, spending energy to keep itself organized, at the small but real price that erasing and re-writing information always carries. A living thing is exactly such a pump, and the arrow favors it, because the growing low-entropy edge keeps supplying the gradient that organized, self-maintaining, reproducing patterns can ride. This is the substrate's version of dissipative adaptation, and the same combinatorial sparseness that makes complex structures rare makes life a thin, threshold-gated tail of all the matter there is, which the survey across the substrate confirms rather than assumes.

Mind is what a self becomes when its memory is the bulk. The content-addressable, logarithmic-depth database of the computation chapter, read from the inside, is a mindscape, where recall is lookup by content, recognition is falling into a stored basin, reasoning is routing through the web, and intuition is a shortcut through the depth. We do not repeat that machinery here, only note that the inner architecture the later chapters open is this same web, felt rather than measured.

Intelligence can be said precisely, and saying it precisely keeps it from being confused with consciousness, which is a different thing entirely. A \emph{problem} is a gap between a desired state and the current state, written $P = D - C$, and when $P$ is zero there is nothing to solve. Intelligence is problem-solving capability, the capacity to drive $P$ to zero, and it is a matter of how well a self can route through its own web toward a target. Consciousness is not that. Consciousness is integrated experience, the felt interior of a self, the fact that there is something it is like to be it, and the two come apart, a system can be highly intelligent with little integrated feeling, or deeply feeling with little problem-solving power.

\begin{principle}[Consciousness as integration]
Consciousness is the degree to which a self's experience is bound into one whole. It is graded, rising with integration, and it rests on the one premise the experiments do not test, that the substrate feels. Granting that premise, the combination problem, how many small feelings make one, is answered structurally, by binding, which is the same integration that makes a self a self.
\end{principle}

Freedom is the last rung, and on a reversible deterministic substrate it has to be argued for rather than assumed. The argument has four strands and concedes the one it must. A self is not overridden by the dynamics, it \emph{is} the deciding part of the dynamics, so at its own level it is the author of its acts, the compatibilist's point. No self holds a complete model of itself, so it cannot predict its own next move, and the future is genuinely open from the inside even where it is fixed from outside. The growing edge keeps making real novelty, so the world is not a finished tape being replayed. And freedom is also a matter of degree a self can gain, by aligning its own motion with the arrow, the contemplative's liberation. What is ruled out is libertarian free will, an uncaused will reaching in from nowhere to break the causal closure, which the base does not permit and the honest account does not smuggle back in.
